% Part: incompleteness
% Chapter: incompleteness-provability
% Section: first-incompleteness-thm

\documentclass[../../../include/open-logic-section]{subfiles}

\begin{document}

\olfileid{inc}{inp}{1in}

\olsection{مبرهنة عدم الاكتمال الأولى}

يمكننا الآن وصف برهان غودل الأصلي لمبرهنة عدم الاكتمال الأولى. لتكن
$\Th{T}$ أي نظرية مؤكسمة حسابيًا في لغة تمتد لغة الحساب، ولتشتمل
$\Th{T}$ على بديهيات~$\Th{Q}$. وهذا يعني، على وجه الخصوص، أن
$\Th{T}$ تمثل الدوال والعلاقات القابلة للحساب.

سبق أن حاججنا بأن العلاقة $\Prf[\Th{T}](x,y)$ قابلة للحساب، متى
اخترنا ترميزًا معقولًا للصيغ والبراهين بأعداد، وأن
$\Prf[\Th{T}](x,y)$ تصدق إذا وفقط إذا كان $x$ عدد غودل
لـ!!a{derivation} للـ!!{formula} ذات عدد غودل~$y$ في~$\Th{T}$. بل
استطاع غودل أن يبيّن، للنظرية المعينة التي كانت في ذهنه، أن هذه
العلاقة عودية بدائية، مستعملًا قائمة الدوال والعلاقات الخمس والأربعين
في بحثه. والعلاقة الخامسة والأربعون، $x B y$، ليست إلا
$\Prf[\Th{T}](x,y)$ بالنسبة إلى اختياره المعين لـ$\Th{T}$. وتذكّر
أننا نستعمل اليوم عبارة «عودية بدائية» حيث يستعمل غودل كلمة «عودية»
في بحثه.

لما كانت $\Prf[\Th{T}](x,y)$ قابلة للحساب، فهي قابلة للتمثيل في
$\Th{T}$. وسنستعمل $\OPrf[\Th{T}](x,y)$ للإشارة إلى الصيغة التي
تمثلها. ولتكن $\OProv[\Th{T}](y)$ الصيغة
$\lexists[x][\OPrf[\Th{T}](x,y)]$. وهذه تصف العلاقة السادسة
والأربعين، $\fn{Bew}(y)$، في قائمة غودل. وكما يلاحظ غودل، فهذه هي
العلاقة الوحيدة التي «لا يمكن الجزم بأنها عودية». ولعله كان يعني ما
يأتي: لا يتضح من التعريف أنها قابلة للحساب، وقد بيّنت التطورات اللاحقة
أنها ليست كذلك بالفعل.

لتكن $\Th{T}$ نظرية !!{axiomatizable} تحتوي~$\Th{Q}$. عندئذ تكون
$\Prf[\Th{T}](x,y)$ قابلة للقرار، ومن ثم تمثلها في~$\Th{Q}$
!!a{formula}~$\OPrf[\Th{T}](x,y)$. ولتكن $\OProv[\Th{T}](y)$ الصيغة
التي وصفناها أعلاه. وبحسب لمّة النقطة الثابتة، توجد صيغة
$!G_\Th{T}$ بحيث !!{derive}s $\Th{Q}$ (ومن ثم $\Th{T}$)
\begin{equation}
\ollabel{eqn:qpf}
!G_\Th{T} \liff \lnot \OProv[\Th{T}](\gn{!G_\Th{T}}).
\end{equation}
لاحظ أن $!G_\Th{T}$ تقول، في جوهرها، «إن $!G_\Th{T}$ غير
!!{derivable} في~$\Th{T}$».

\begin{lem}\ollabel{lem:cons-G-unprov}
إذا كانت $\Th{T}$ نظرية متسقة و!!{axiomatizable} تمتدّ~$\Th{Q}$،
فإن $\Th{T}\Proves/ !G_\Th{T}$.
\end{lem}

\begin{proof}
افترض أن النظام~$\Th{T}$ !!{derive}s $!G_\Th{T}$. عندئذ يوجد
\emph{بالفعل} !!a{derivation}، ولذلك تصدق، لعدد ما~$m$، العلاقة
$\Prf[\Th{T}](m,\Gn{!G_\Th{T}})$. لكن $\Th{Q}$ !!{derive}s عندئذ
الجملة $\OPrf[\Th{T}](\num m,\gn{!G_\Th{T}})$. ومن ثم !!{derive}s
$\Th{Q}$ الصيغة
$\lexists[x][\OPrf[\Th{T}](x,\gn{!G_\Th{T}})]$، وهي، بالتعريف،
$\OProv[\Th{T}](\gn{!G_\Th{T}})$. وبحسب \olref{eqn:qpf}،
!!{derive}s $\Th{Q}$ الصيغة $\lnot !G_\Th{T}$، ولما كانت $\Th{T}$
تمتدّ $\Th{Q}$ فإن $\Th{T}$ !!{derive}s إياها أيضًا. لقد بيّنا أنه
إذا كان النظام~$\Th{T}$ !!{derive}s $!G_\Th{T}$، فإنه !!{derive}s
أيضًا $\lnot !G_\Th{T}$، فيكون غير متسق.
\end{proof}

\begin{defn}
\ollabel{thm:oconsis-q}
تكون النظرية $\Th{T}$ \emph{متسقة-$\omega$} إذا تحقق الآتي: إذا
كانت $\lexists[x][!A(x)]$ أي جملة، وكان النظام~$\Th{T}$ !!{derive}s
$\lnot !A(\num 0)$، و$\lnot !A(\num 1)$، و$\lnot !A(\num 2)$،
و\dots، فإن $\Th{T}$ لا تثبت $\lexists[x][!A(x)]$.
\end{defn}

لاحظ أن كل نظرية متسقة-$\omega$ متسقة أيضًا. وينتج هذا مباشرة من أنه
إذا كانت $\Th{T}$ غير متسقة، فإن $\Th{T}\Proves !A$ لكل~$!A$. وعلى وجه
الخصوص، إذا كانت $\Th{T}$ غير متسقة، فإنها !!{derive}s
كلًا من $\lnot !A(\num n)$ لكل~$n$ و
$\lexists[x][!A(x)]$. ولذلك، إذا كانت $\Th{T}$ غير متسقة، فهي غير
متسقة-$\omega$. وبعكس النقيض، إذا كانت $\Th{T}$ متسقة-$\omega$، فلا
بد أن تكون متسقة.

\begin{lem}\ollabel{lem:omega-cons-G-unref}
إذا كانت $\Th{T}$ نظرية متسقة-$\omega$ و!!{axiomatizable}
تمتدّ~$\Th{Q}$، فإن $\Th{T}\Proves/\lnot !G_\Th{T}$.
\end{lem}

\begin{proof}
سنبيّن أنه إذا كان النظام~$\Th{T}$ !!{derive}s $\lnot !G_\Th{T}$،
فهو غير متسق-$\omega$. افترض أن النظام~$\Th{T}$ !!{derive}s
$\lnot !G_\Th{T}$. فإذا كانت $\Th{T}$ غير متسقة، كانت غير
متسقة-$\omega$ وانتهينا. وإلا كانت $\Th{T}$ متسقة، فلا يكون في وسعها !!{derive}
$!G_\Th{T}$ بحسب \olref{lem:cons-G-unprov}. ولما لم يوجد
!!{derivation} لـ$!G_\Th{T}$ في~$\Th{T}$، فإن $\Th{Q}$ !!{derive}s
\[
\lnot \OPrf[\Th{T}](\num 0, \gn{!G_\Th{T}}), \lnot \OPrf[\Th{T}](\num 1,
\gn{!G_\Th{T}}), \lnot \OPrf[\Th{T}](\num 2, \gn{!G_\Th{T}}), \dots
\]
وكذلك تفعل~$\Th{T}$. ومن جهة أخرى، تكافئ $\lnot !G_\Th{T}$، بحسب
\olref{eqn:qpf}، الصيغة
$\lexists[x][\OPrf[\Th{T}](x,\gn{!G_\Th{T}})]$. ومن ثم تكون
$\Th{T}$ غير متسقة-$\omega$.
\end{proof}

\begin{prob}
  كل نظرية متسقة-$\omega$ متسقة. بيّن أن العكس لا يصح؛ أي إن ثمة
  نظريات متسقة وغير متسقة-$\omega$. وافعل ذلك ببيان أن
  $\Th{Q}\cup\{\lnot !G_\Th{Q}\}$ متسقة وغير متسقة-$\omega$.
\end{prob}

\begin{thm}
\ollabel{thm:first-incompleteness} لتكن $\Th{T}$ أي نظرية
متسقة-$\omega$ و!!{axiomatizable} تمتدّ~$\Th{Q}$. عندئذ لا تكون
$\Th{T}$ تامة.
\end{thm}

\begin{proof}
  إذا كانت $\Th{T}$ متسقة-$\omega$، فهي متسقة، ومن ثم
  $\Th{T}\Proves/ !G_\Th{T}$ بحسب \olref{lem:cons-G-unprov}.
  وبحسب \olref{lem:omega-cons-G-unref}،
  $\Th{T}\Proves/\lnot !G_\Th{T}$. وهذا يعني أن $\Th{T}$ غير تامة،
  لأنها لا !!{derive}s $!G_\Th{T}$ ولا $\lnot !G_\Th{T}$.
\end{proof}

\end{document}
